\documentclass[11pt]{article}

% Shared notes settings for Elements of AIML lectures (UPES).
% Keep this file free of \documentclass to allow reuse via \input.

\usepackage[utf8]{inputenc}
\usepackage[T1]{fontenc}
\usepackage[a4paper,margin=1in]{geometry}
\usepackage{hyperref}
\usepackage{xurl}
\usepackage{booktabs}
\usepackage{amsmath}
\usepackage{amssymb}
\usepackage{listings}
\usepackage{xcolor}
\usepackage{enumitem}
\usepackage{parskip}

% Code listings (general-purpose).
\lstset{
  basicstyle=\ttfamily\small,
  keywordstyle=\color{blue},
  commentstyle=\color{gray},
  stringstyle=\color{teal},
  showstringspaces=false,
  columns=fullflexible,
  frame=single,
  framerule=0.2pt,
  breaklines=true
}

\setlist[itemize]{topsep=0.3em,itemsep=0.2em}
\setlist[enumerate]{topsep=0.3em,itemsep=0.2em}

% Simple per-section source line.
\newcommand{\notesources}[1]{\par\noindent{\small\color{gray}\textit{Sources:} \sloppy #1\par}}

% Unit-wise source strings for this subject (used by slides/notes).
% Path is relative to the lecture `latex/` folder (the compilation CWD).
\IfFileExists{../../../_shared/aiml-sources.tex}{%
  % Source strings for Elements of AIML (used by slides + notes).

\newcommand{\AIMLRepoURL}{https://github.com/tali7c/Elements-of-AIML}

\newcommand{\AIMLLabSources}{%
Provost and Fawcett, \textit{Data Science for Business}.;%
McKinney, \textit{Python for Data Analysis}.;%
WEKA Documentation (Waikato Environment for Knowledge Analysis), accessed 2026-02-28.;%
Matplotlib and Seaborn documentation, accessed 2026-02-28.%
}

\newcommand{\AIMLUnitISources}{%
Rich and Knight, \textit{Artificial Intelligence}, McGraw Hill, 2017.;%
Russell and Norvig, \textit{Artificial Intelligence: A Modern Approach} (4e), Ch. 1--2 (Introduction, Intelligent Agents).;%
Poole and Mackworth, \textit{Artificial Intelligence: Foundations of Computational Agents} (2e), selected sections.%
}

\newcommand{\AIMLUnitIISources}{%
Russell and Norvig, \textit{Artificial Intelligence: A Modern Approach} (4e), Ch. 7--9 (Logical Agents, First-Order Logic, Inference).;%
Huth and Ryan, \textit{Logic in Computer Science} (2e), selected sections on propositional and first-order logic.;%
Genesereth and Nilsson, \textit{Logical Foundations of Artificial Intelligence}, selected chapters.%
}

\newcommand{\AIMLUnitIIISources}{%
Mueller and Massaron, \textit{Machine Learning for Dummies}, For Dummies, 2016.;%
Mitchell, \textit{Machine Learning}, selected chapters on learning basics and evaluation.;%
Scikit-learn documentation, accessed 2026-02-28.%
}

\newcommand{\AIMLUnitIVSources}{%
Mueller and Massaron, \textit{Machine Learning for Dummies}, For Dummies, 2016.;%
James, Witten, Hastie, and Tibshirani, \textit{An Introduction to Statistical Learning} (2e), selected chapters.;%
Scikit-learn documentation, accessed 2026-02-28.%
}

\newcommand{\AIMLUnitVSources}{%
NIST, \textit{AI Risk Management Framework (AI RMF 1.0)}, 2023.;%
Barocas, Hardt, and Narayanan, \textit{Fairness and Machine Learning} (online book), selected sections.;%
Knaflic, \textit{Storytelling with Data}, selected chapters on communicating results.%
}
%
}{%
  % Faculty copies live under `private/` so their compile CWD is different.
  \IfFileExists{../../../../public/_shared/aiml-sources.tex}{%
    % Source strings for Elements of AIML (used by slides + notes).

\newcommand{\AIMLRepoURL}{https://github.com/tali7c/Elements-of-AIML}

\newcommand{\AIMLLabSources}{%
Provost and Fawcett, \textit{Data Science for Business}.;%
McKinney, \textit{Python for Data Analysis}.;%
WEKA Documentation (Waikato Environment for Knowledge Analysis), accessed 2026-02-28.;%
Matplotlib and Seaborn documentation, accessed 2026-02-28.%
}

\newcommand{\AIMLUnitISources}{%
Rich and Knight, \textit{Artificial Intelligence}, McGraw Hill, 2017.;%
Russell and Norvig, \textit{Artificial Intelligence: A Modern Approach} (4e), Ch. 1--2 (Introduction, Intelligent Agents).;%
Poole and Mackworth, \textit{Artificial Intelligence: Foundations of Computational Agents} (2e), selected sections.%
}

\newcommand{\AIMLUnitIISources}{%
Russell and Norvig, \textit{Artificial Intelligence: A Modern Approach} (4e), Ch. 7--9 (Logical Agents, First-Order Logic, Inference).;%
Huth and Ryan, \textit{Logic in Computer Science} (2e), selected sections on propositional and first-order logic.;%
Genesereth and Nilsson, \textit{Logical Foundations of Artificial Intelligence}, selected chapters.%
}

\newcommand{\AIMLUnitIIISources}{%
Mueller and Massaron, \textit{Machine Learning for Dummies}, For Dummies, 2016.;%
Mitchell, \textit{Machine Learning}, selected chapters on learning basics and evaluation.;%
Scikit-learn documentation, accessed 2026-02-28.%
}

\newcommand{\AIMLUnitIVSources}{%
Mueller and Massaron, \textit{Machine Learning for Dummies}, For Dummies, 2016.;%
James, Witten, Hastie, and Tibshirani, \textit{An Introduction to Statistical Learning} (2e), selected chapters.;%
Scikit-learn documentation, accessed 2026-02-28.%
}

\newcommand{\AIMLUnitVSources}{%
NIST, \textit{AI Risk Management Framework (AI RMF 1.0)}, 2023.;%
Barocas, Hardt, and Narayanan, \textit{Fairness and Machine Learning} (online book), selected sections.;%
Knaflic, \textit{Storytelling with Data}, selected chapters on communicating results.%
}
%
  }{%
    % Source strings for Elements of AIML (used by slides + notes).

\newcommand{\AIMLRepoURL}{https://github.com/tali7c/Elements-of-AIML}

\newcommand{\AIMLLabSources}{%
Provost and Fawcett, \textit{Data Science for Business}.;%
McKinney, \textit{Python for Data Analysis}.;%
WEKA Documentation (Waikato Environment for Knowledge Analysis), accessed 2026-02-28.;%
Matplotlib and Seaborn documentation, accessed 2026-02-28.%
}

\newcommand{\AIMLUnitISources}{%
Rich and Knight, \textit{Artificial Intelligence}, McGraw Hill, 2017.;%
Russell and Norvig, \textit{Artificial Intelligence: A Modern Approach} (4e), Ch. 1--2 (Introduction, Intelligent Agents).;%
Poole and Mackworth, \textit{Artificial Intelligence: Foundations of Computational Agents} (2e), selected sections.%
}

\newcommand{\AIMLUnitIISources}{%
Russell and Norvig, \textit{Artificial Intelligence: A Modern Approach} (4e), Ch. 7--9 (Logical Agents, First-Order Logic, Inference).;%
Huth and Ryan, \textit{Logic in Computer Science} (2e), selected sections on propositional and first-order logic.;%
Genesereth and Nilsson, \textit{Logical Foundations of Artificial Intelligence}, selected chapters.%
}

\newcommand{\AIMLUnitIIISources}{%
Mueller and Massaron, \textit{Machine Learning for Dummies}, For Dummies, 2016.;%
Mitchell, \textit{Machine Learning}, selected chapters on learning basics and evaluation.;%
Scikit-learn documentation, accessed 2026-02-28.%
}

\newcommand{\AIMLUnitIVSources}{%
Mueller and Massaron, \textit{Machine Learning for Dummies}, For Dummies, 2016.;%
James, Witten, Hastie, and Tibshirani, \textit{An Introduction to Statistical Learning} (2e), selected chapters.;%
Scikit-learn documentation, accessed 2026-02-28.%
}

\newcommand{\AIMLUnitVSources}{%
NIST, \textit{AI Risk Management Framework (AI RMF 1.0)}, 2023.;%
Barocas, Hardt, and Narayanan, \textit{Fairness and Machine Learning} (online book), selected sections.;%
Knaflic, \textit{Storytelling with Data}, selected chapters on communicating results.%
}
%
  }%
}


\title{Elements of AIML\\Unit 02 -- Lecture 03 Notes\\Clause Form, Resolution \& Unification}
\author{Tofik Ali}
\date{\today}

\begin{document}
\maketitle
\tableofcontents

\section{Lecture Overview}
This lecture introduces \textbf{resolution}, a widely used inference method for automated reasoning, and the supporting ideas:
\begin{itemize}[leftmargin=*]
  \item converting sentences to \textbf{clause form} (CNF),
  \item proving entailment by \textbf{refutation} ($KB \land \lnot \alpha$ unsatisfiable),
  \item \textbf{unification} to extend resolution from propositional logic to first-order logic.
\end{itemize}
\notesources{\AIMLUnitIISources}

\section{How to use these notes (self-study)}
Resolution becomes easy only after you are comfortable with two routines:
(1) converting sentences to CNF (clause sets), and (2) performing resolution steps carefully until no new clauses appear.
For FOL resolution, add a third routine: unification (computing substitutions).
While studying, always keep a paper ``clause set'' and write each new resolvent explicitly.

\section{Learning Outcomes}
\begin{itemize}[leftmargin=*]
  \item Convert simple propositional statements to CNF.
  \item Use resolution to show entailment in propositional logic for small KBs.
  \item Explain the main extra steps needed for FOL resolution.
  \item Define unification and compute simple MGUs.
\end{itemize}

\section{Keywords}
\begin{itemize}[leftmargin=*]
  \item \textbf{CNF:} conjunctive normal form (AND of OR clauses).
  \item \textbf{Literal:} an atom or its negation.
  \item \textbf{Clause:} a disjunction of literals.
  \item \textbf{Resolution:} rule of inference on clauses (refutation proof method).
  \item \textbf{Refutation:} prove $KB \models \alpha$ by showing $KB \land \lnot\alpha$ is unsatisfiable.
  \item \textbf{Unification:} find a substitution to make expressions identical.
  \item \textbf{MGU:} most general unifier (least restrictive substitution).
\end{itemize}

\section{Abbreviations and symbols}
\begin{itemize}[leftmargin=*]
  \item \textbf{PL}: propositional logic.
  \item \textbf{FOL}: first-order logic (predicate logic).
  \item $KB$: knowledge base (set of sentences/facts/rules).
  \item $\alpha$: query/goal sentence.
  \item $\models$: semantic entailment (truth in \emph{all} models).
  \item $\vdash$: syntactic derivation/provability in a proof system.
  \item $\lnot,\land,\lor,\rightarrow,\leftrightarrow$: NOT, AND, OR, IMPLIES, IFF.
  \item \textbf{CNF}: conjunctive normal form; \textbf{DNF}: disjunctive normal form.
  \item $\Box$: empty clause (contradiction).
  \item $\theta$: a substitution (mapping variables $\rightarrow$ terms); \textbf{MGU}: most general unifier.
  \item \textbf{Skolem constant/function}: new symbols introduced to eliminate $\exists$ in clause form.
\end{itemize}

\section{Refutation view of entailment}
In logic for AI we distinguish:
\begin{itemize}[leftmargin=*]
  \item \textbf{Semantics (meaning):} what is true in the world(s) described by our sentences.
  \item \textbf{Syntax (proof):} how we can mechanically derive new sentences from existing ones.
\end{itemize}

Two important relations:
\begin{itemize}[leftmargin=*]
  \item \textbf{Entailment} $KB \models \alpha$: in \emph{every} model where all sentences in $KB$ are true, $\alpha$ is also true.
  \item \textbf{Derivation} $KB \vdash \alpha$: $\alpha$ can be proven from $KB$ using a chosen proof procedure (rules of inference).
\end{itemize}

If a proof system is \textbf{sound}, then $KB \vdash \alpha \Rightarrow KB \models \alpha$ (no false conclusions).
If it is \textbf{complete}, then $KB \models \alpha \Rightarrow KB \vdash \alpha$ (it can prove everything that is true).

The central equivalence used by resolution:
\[
  KB \models \alpha \;\;\Leftrightarrow\;\; KB \land \lnot \alpha \text{ is unsatisfiable}
\]
So, to prove entailment:
\begin{enumerate}[leftmargin=*]
  \item Take all sentences in $KB$ and the negation of the query $\lnot\alpha$.
  \item Convert them into CNF (set of clauses).
  \item Apply resolution repeatedly.
  \item If you derive the \textbf{empty clause} ($\Box$), a contradiction is reached, so $KB \models \alpha$.
\end{enumerate}

\section{Clause form (CNF) in propositional logic}
\subsection{What exactly is a clause set?}
In resolution we usually treat CNF as a \textbf{set of clauses}:
\[
  (C_1) \land (C_2) \land \dots \land (C_m)
\]
where each $C_i$ is a disjunction of literals.
For automated reasoning it is convenient to store CNF as:
\[
  \{C_1, C_2, \dots, C_m\}
\]
i.e., a \emph{set} (or list) of clauses.

\subsection{CNF conversion recipe (PL)}
\begin{enumerate}[leftmargin=*]
  \item Eliminate $\leftrightarrow$ and $\rightarrow$.
  \item Push negations inward (De Morgan + double negation).
  \item Distribute $\lor$ over $\land$ to obtain an AND of ORs.
  \item Flatten nested AND/OR groups.
\end{enumerate}

\subsection{Mini example}
Convert $(p \rightarrow q) \land p$:
\begin{enumerate}[leftmargin=*]
  \item $(\lnot p \lor q) \land p$
  \item CNF clauses: $(\lnot p \lor q)$ and $(p)$
\end{enumerate}

\subsection{Worked example (distribution is required)}
Convert $p \rightarrow (q \land r)$:
\begin{enumerate}[leftmargin=*]
  \item Eliminate implication:
  \[
    p \rightarrow (q \land r)\;\;\equiv\;\; \lnot p \lor (q \land r)
  \]
  \item Distribute $\lor$ over $\land$:
  \[
    \lnot p \lor (q \land r)\;\;\equiv\;\;(\lnot p \lor q) \land (\lnot p \lor r)
  \]
  \item Clauses: $(\lnot p \lor q)$ and $(\lnot p \lor r)$.
\end{enumerate}

\subsection{Worked example (biconditional)}
Convert $p \leftrightarrow q$:
\begin{enumerate}[leftmargin=*]
  \item Replace IFF:
  \[
    (p \leftrightarrow q) \equiv (p \rightarrow q) \land (q \rightarrow p)
  \]
  \item Eliminate implications:
  \[
    (\lnot p \lor q) \land (\lnot q \lor p)
  \]
  \item Already in CNF (two clauses).
\end{enumerate}

\subsection{Useful CNF simplifications (optional but practical)}
After conversion we often simplify the clause set:
\begin{itemize}[leftmargin=*]
  \item Remove duplicate literals inside a clause: $(p \lor p \lor q) \Rightarrow (p \lor q)$.
  \item Remove clauses that are always true (tautological): $(p \lor \lnot p \lor q)$ can be dropped.
  \item Use \textbf{unit clauses} like $(p)$ effectively (unit resolution).
\end{itemize}

\section{Resolution in propositional logic}
\subsection{Resolution rule}
Given two clauses $(A \lor p)$ and $(B \lor \lnot p)$, resolve on $p$ to infer $(A \lor B)$.

\subsection{Why resolution is useful}
\begin{itemize}[leftmargin=*]
  \item It is \textbf{sound}: any resolvent derived is logically implied by the parents.
  \item For propositional logic, it is also \textbf{complete} for refutation: if the clause set is unsatisfiable, resolution can derive $\Box$.
  \item It reduces reasoning to a small number of mechanical steps, which is why SAT solvers and automated provers are built around similar ideas.
\end{itemize}

\subsection{Worked proof: $(p\rightarrow q), p \models q$}
Let $KB=\{p\rightarrow q,\; p\}$ and query $\alpha=q$.
\begin{enumerate}[leftmargin=*]
  \item Add $\lnot q$ (negated query).
  \item Convert to CNF clauses:
  \[
    (\lnot p \lor q),\; (p),\; (\lnot q)
  \]
  \item Resolve $(\lnot p \lor q)$ with $(p)$ to get $(q)$.
  \item Resolve $(q)$ with $(\lnot q)$ to get $\Box$ (empty clause).
\end{enumerate}
Hence $KB \models q$.

\subsection{Worked proof: $(p\rightarrow q),(q\rightarrow r), p \models r$}
Let $KB=\{p\rightarrow q,\; q\rightarrow r,\; p\}$ and query $\alpha=r$.
\begin{enumerate}[leftmargin=*]
  \item Add $\lnot r$.
  \item CNF clauses:
  \[
    C_1:(\lnot p \lor q),\quad
    C_2:(\lnot q \lor r),\quad
    C_3:(p),\quad
    C_4:(\lnot r)
  \]
  \item Resolve $C_1$ with $C_3$ on $p$ to get $C_5:(q)$.
  \item Resolve $C_2$ with $C_5$ on $q$ to get $C_6:(r)$.
  \item Resolve $C_6$ with $C_4$ on $r$ to derive $\Box$.
\end{enumerate}
Therefore $KB \models r$.

\subsection{Example: detecting inconsistency (unsatisfiable KB)}
Suppose we have:
\[
  (p \lor q),\quad (\lnot p),\quad (\lnot q)
\]
Resolve $(p \lor q)$ with $(\lnot p)$ to get $(q)$, then resolve $(q)$ with $(\lnot q)$ to get $\Box$.
So the set is \textbf{unsatisfiable} (inconsistent).

\subsection{Algorithmic view (how an automated prover uses it)}
An automated prover typically:
\begin{itemize}[leftmargin=*]
  \item maintains a growing set of clauses,
  \item repeatedly chooses pairs of clauses with complementary literals,
  \item adds resolvents until either $\Box$ is derived or no new clauses can be produced.
\end{itemize}
In practice, \textbf{search strategy} matters (unit preference, set-of-support, subsumption), but the core logical rule stays the same.

\section{Resolution in first-order logic (high-level)}
FOL resolution is similar, but clause form conversion can include additional steps (presented at a conceptual level here):
\begin{itemize}[leftmargin=*]
  \item remove implications,
  \item move negations inward,
  \item standardize variables (avoid clashes),
  \item handle quantifiers (often by Skolemization for $\exists$),
  \item drop universal quantifiers (implicit in clauses),
  \item produce a set of clauses.
\end{itemize}
Then resolution uses unification to match literals.

\subsection{Clause form conversion in FOL (more detail)}
To convert a FOL sentence to clause form for resolution, a common pipeline is:
\begin{enumerate}[leftmargin=*]
  \item Eliminate $\leftrightarrow$ and $\rightarrow$.
  \item Push $\lnot$ inward so that it applies only to atomic predicates.
  \item Rename variables so that different sentences do not accidentally share variable names (\textbf{standardize apart}).
  \item Move quantifiers to the front (prenex form) if needed.
  \item Eliminate $\exists$ using \textbf{Skolemization}:
  \begin{itemize}[leftmargin=*]
    \item if $\exists y$ is inside the scope of $\forall x$, replace $y$ by a new function $f(x)$;
    \item if $\exists y$ has no surrounding $\forall$, replace $y$ by a new constant $k$.
  \end{itemize}
  \item Drop $\forall$ (universals are implicit in clause variables).
  \item Distribute $\lor$ over $\land$, then split into clauses.
\end{enumerate}

\subsection{Worked clause form example (Skolemization)}
Sentence:
\[
  \forall x\; (Student(x) \rightarrow \exists y\; Enrolled(x,y))
  \quad\text{``Every student is enrolled in some course.''}
\]
Steps:
\begin{enumerate}[leftmargin=*]
  \item Remove implication:
  \[
    \forall x\; (\lnot Student(x) \lor \exists y\; Enrolled(x,y))
  \]
  \item Pull $\exists$ out (since $y$ does not appear in $\lnot Student(x)$):
  \[
    \forall x\; \exists y\; (\lnot Student(x) \lor Enrolled(x,y))
  \]
  \item Skolemize $y$ with a new function $f(x)$:
  \[
    \forall x\; (\lnot Student(x) \lor Enrolled(x,f(x)))
  \]
  \item Drop $\forall$ and read as one clause:
  \[
    \lnot Student(x) \lor Enrolled(x,f(x))
  \]
\end{enumerate}
The Skolem function $f(x)$ should be read as ``the course chosen for $x$'' (it may differ for different $x$).

\subsection{Classic example: Socrates is mortal}
Knowledge:
\[
  \forall x\; (Human(x) \rightarrow Mortal(x)),\qquad Human(Socrates)
\]
Query: $Mortal(Socrates)$.

Convert to clause-like form:
\begin{itemize}[leftmargin=*]
  \item $Human(x) \rightarrow Mortal(x)$ becomes $\lnot Human(x) \lor Mortal(x)$.
  \item Add $\lnot Mortal(Socrates)$ (negated query).
\end{itemize}
Clauses:
\[
  (\lnot Human(x) \lor Mortal(x)),\quad Human(Socrates),\quad \lnot Mortal(Socrates)
\]
Resolve $(\lnot Human(x) \lor Mortal(x))$ with $Human(Socrates)$ using substitution $\{x/Socrates\}$ to get $Mortal(Socrates)$, then resolve with $\lnot Mortal(Socrates)$ to reach $\Box$.

\section{Unification}
\subsection{Definition}
Unification finds a substitution $\theta$ such that $E_1\theta = E_2\theta$.

\subsection{Substitutions (what $\theta$ really is)}
A substitution $\theta$ is a finite set of bindings like:
\[
  \theta = \{x/t,\; y/u,\; z/v\}
\]
meaning ``replace $x$ by term $t$, $y$ by term $u$, etc.''.
Applying $\theta$ to an expression is written $E\theta$.

\subsection{Most general unifier (MGU)}
There may be many unifiers; we prefer the \textbf{most general} one:
\begin{itemize}[leftmargin=*]
  \item $\theta$ is more general than $\sigma$ if $\sigma$ can be obtained by adding extra substitutions after $\theta$.
  \item MGU keeps as many variables free as possible (least restrictive).
\end{itemize}

\subsection{Unification algorithm (high-level idea)}
Unification can be seen as repeatedly simplifying a set of equations.
One simple conceptual view:
\begin{lstlisting}
UNIFY(E1, E2):
  if E1 == E2: return {}
  if E1 is a variable x: return {x / E2} (if occurs-check passes)
  if E2 is a variable x: return {x / E1} (if occurs-check passes)
  if E1 = f(a1,...,an) and E2 = f(b1,...,bn):
       unify each pair (ai, bi) and combine substitutions
  else: fail
\end{lstlisting}
In real provers, the implementation is more careful, but the cases above explain the intuition.

\subsection{Examples}
\begin{itemize}[leftmargin=*]
  \item Unify $Likes(x, AI)$ with $Likes(Ali, AI)$:
  \[
    \theta = \{x/Ali\}
  \]
  \item Unify $Parent(x, y)$ with $Parent(Ali, Sara)$:
  \[
     \theta = \{x/Ali,\; y/Sara\}
  \]
  \item Unify $Knows(x, f(y), y)$ with $Knows(Ali, f(z), z)$:
  \[
    \theta = \{x/Ali,\; y/z\}
  \]
\end{itemize}

\subsection{Occurs-check intuition}
Unifying $x$ with $f(x)$ would require an infinite structure. So a correct unification algorithm rejects such cases (the occurs-check).

\section{Practice (with answers)}

\subsection*{P1. Convert to CNF: $(p \rightarrow q) \land (q \rightarrow r)$.}
\textbf{Answer:} $(\lnot p \lor q) \land (\lnot q \lor r)$.

\subsection*{P2. Show by resolution: $(p \rightarrow q), (q \rightarrow r), p \models r$.}
\textbf{Answer (details):} Add $\lnot r$. CNF clauses: $(\lnot p \lor q)$, $(\lnot q \lor r)$, $(p)$, $(\lnot r)$. Resolve $(\lnot p \lor q)$ with $(p)$ to get $(q)$. Resolve $(\lnot q \lor r)$ with $(q)$ to get $(r)$. Resolve $(r)$ with $(\lnot r)$ to get $\Box$.

\subsection*{P3. Find an MGU for $Knows(x, f(y))$ and $Knows(Ali, f(z))$.}
\textbf{Answer:} $\{x/Ali,\; y/z\}$ (or equivalently $\{x/Ali,\; z/y\}$ depending on variable choice).

\subsection*{P4. Convert to CNF: $p \rightarrow (q \land r)$.}
\textbf{Answer:} $\lnot p \lor (q \land r) \equiv (\lnot p \lor q) \land (\lnot p \lor r)$.

\subsection*{P5. Convert to clause form: $\forall x\; (P(x) \rightarrow \exists y\; Q(x,y))$.}
\textbf{Answer:} $\forall x\; (\lnot P(x) \lor \exists y\; Q(x,y)) \equiv \forall x \exists y (\lnot P(x) \lor Q(x,y))$. Skolemize $y$ as $f(x)$: $\forall x (\lnot P(x) \lor Q(x,f(x)))$. Drop $\forall$: clause $(\lnot P(x) \lor Q(x,f(x)))$.

\subsection*{P6. Give a small inconsistent clause set and show $\Box$ by resolution.}
\textbf{Answer:} $(p \lor q),\;(\lnot p),\;(\lnot q)$. Resolve $(p \lor q)$ with $(\lnot p)$ to get $(q)$, then resolve $(q)$ with $(\lnot q)$ to get $\Box$.

\section{Exit questions (with answers)}

\subsection*{Q0. Entailment vs derivation?}
\textbf{Answer:} $KB \models \alpha$ is a semantic claim about truth in all models; $KB \vdash \alpha$ means there is a proof of $\alpha$ from $KB$ in a chosen proof system.

\subsection*{Q1. Why do we convert to CNF before resolution?}
\textbf{Answer:} the resolution rule operates on clauses (disjunctions of literals). CNF provides a standard clause representation enabling systematic inference.

\subsection*{Q2. State the resolution rule for propositional logic.}
\textbf{Answer:} from $(A \lor p)$ and $(B \lor \lnot p)$ infer $(A \lor B)$.

\subsection*{Q3. Explain unification with an example.}
\textbf{Answer:} unification finds a substitution to make expressions identical; e.g., unify $Likes(x,AI)$ with $Likes(Ali,AI)$ using $\{x/Ali\}$.

\subsection*{Q4. Why does the occurs-check matter?}
\textbf{Answer:} it prevents invalid substitutions like $x=f(x)$ that would imply infinite terms and break correctness.

\subsection*{Q5. What does the empty clause $\Box$ mean in resolution?}
\textbf{Answer:} it represents an explicit contradiction: there is no possible truth assignment/model that can satisfy all derived clauses, so the clause set is unsatisfiable.

\subsection*{Q6. Why is Skolemization needed in FOL clause form?}
\textbf{Answer:} resolution works on universally quantified clauses. Skolemization eliminates existential quantifiers by introducing new constants/functions that preserve satisfiability, allowing us to represent ``there exists'' choices inside clauses.

\section{Common pitfalls (and how to avoid them)}
\begin{itemize}[leftmargin=*]
  \item \textbf{Forgetting the refutation setup:} to prove $KB\models \alpha$, you must add $\lnot\alpha$ and try to derive $\Box$.
  \item \textbf{CNF conversion mistakes:} especially distributing $\lor$ over $\land$ and pushing negations inward correctly.
  \item \textbf{Variable capture/clashes in FOL:} standardize variables apart before Skolemization and clause extraction.
  \item \textbf{Treating Skolem symbols as ``real'' names:} they are placeholders introduced to remove $\exists$ while preserving satisfiability.
  \item \textbf{Unification errors:} apply substitutions consistently and remember the occurs-check intuition ($x$ cannot unify with $f(x)$).
\end{itemize}

\section{Summary}
Resolution provides an automated reasoning method by reducing entailment to unsatisfiability, using CNF clauses and (for FOL) unification.
\notesources{\AIMLUnitIISources}

\end{document}
